% Shared Overleaf preamble (Harshibar / jake-yang template).
% In Overleaf: upload this file with the resume-*.tex you want.
% Menu → Compiler → pdfLaTeX. Menu → Main document → the resume-*.tex file.
\documentclass[letterpaper,11pt]{article}

\usepackage{latexsym}
\usepackage[empty]{fullpage}
\usepackage{titlesec}
\usepackage{marvosym}
\usepackage[usenames,dvipsnames]{color}
\usepackage{verbatim}
\usepackage{enumitem}
\usepackage[hidelinks]{hyperref}
\usepackage{fancyhdr}
\usepackage[english]{babel}
\usepackage{tabularx}
\usepackage{iftex}
\usepackage{fontawesome5}

\ifPDFTeX
  \usepackage[T1]{fontenc}
  \usepackage{tgheros}
  \renewcommand*\familydefault{\sfdefault}
  \usepackage[scale=0.90,lf]{FiraMono}
\else
  \usepackage{fontspec}
  \defaultfontfeatures{Ligatures=TeX}
  % Overleaf LuaLaTeX has TeX Gyre Heros; local XeTeX/tectonic may not
  \IfFontExistsTF{TeX Gyre Heros}{
    \setsansfont{TeX Gyre Heros}
  }{
    \IfFontExistsTF{Helvetica Neue}{
      \setsansfont{Helvetica Neue}
    }{
      \setsansfont{Helvetica}
    }
  }
  \IfFontExistsTF{Fira Mono}{
    \setmonofont{Fira Mono}[Scale=0.90]
  }{
    \IfFontExistsTF{Menlo}{
      \setmonofont{Menlo}[Scale=0.90]
    }{
      \setmonofont{Courier New}[Scale=0.90]
    }
  }
  \renewcommand*\familydefault{\sfdefault}
\fi

\definecolor{light-grey}{gray}{0.83}
\definecolor{dark-grey}{gray}{0.3}
\definecolor{text-grey}{gray}{.08}

\DeclareRobustCommand{\ebseries}{\fontseries{eb}\selectfont}
\DeclareTextFontCommand{\texteb}{\ebseries}
\newcommand{\myuline}[1]{\underline{#1}}

\pagestyle{fancy}
\fancyhf{}
\fancyfoot{}
\renewcommand{\headrulewidth}{0pt}
\renewcommand{\footrulewidth}{0pt}

\addtolength{\oddsidemargin}{-0.5in}
\addtolength{\evensidemargin}{0in}
\addtolength{\textwidth}{1in}
\addtolength{\topmargin}{-0.5in}
\addtolength{\textheight}{1.0in}

\urlstyle{same}
\raggedbottom
\raggedright
\setlength{\tabcolsep}{0in}

\titleformat{\section}{
    \bfseries \vspace{1pt} \raggedright \large
}{}{0em}{}[\color{light-grey}{\titlerule[2pt]}\vspace{-6pt}]

\newcommand{\resumeItem}[1]{
  \item\small{
    {#1 \vspace{-3pt}}
  }
}

\newcommand{\resumeSubheading}[4]{
  \vspace{-1pt}\item
    \begin{tabular*}{\textwidth}[t]{l@{\extracolsep{\fill}}r}
      \textbf{#1} & {\color{dark-grey}\small #2}\vspace{1pt}\\
      \textit{#3} & {\color{dark-grey} \small #4}\\
    \end{tabular*}\vspace{-4pt}
}

\newcommand{\resumeProjectHeading}[2]{
    \item
    \begin{tabular*}{\textwidth}{l@{\extracolsep{\fill}}r}
      #1 & {\color{dark-grey} \small #2} \\
    \end{tabular*}\vspace{-4pt}
}

\renewcommand\labelitemii{$\vcenter{\hbox{\tiny$\bullet$}}$}

\newcommand{\resumeSubHeadingListStart}{\begin{itemize}[leftmargin=0in, label={}]}
\newcommand{\resumeSubHeadingListEnd}{\end{itemize}}
\newcommand{\resumeItemListStart}{\begin{itemize}[topsep=1pt, itemsep=1pt, parsep=0pt, partopsep=0pt]}
\newcommand{\resumeItemListEnd}{\end{itemize}\vspace{-2pt}}

\color{text-grey}
